\documentclass[a4paper]{article}
\usepackage{amsfonts}
\usepackage{amsmath}
\usepackage{amssymb}
\usepackage{graphicx}     

\begin{document}
  
\noindent \large Auteur: Stan Van Campenhout \\
\noindent \large Studierichting:  Informatica\\
\noindent \large Oefeningenreeks 5A nr. 2.3\\

\medskip

\normalsize

$ \displaystyle\int\limits_{-1}^2\left(x\cdot\lfloor x\rfloor\right) dx = \displaystyle\int\limits_{-1}^0\left(x\cdot\lfloor x\rfloor\right) dx 
+ \displaystyle\int\limits_0^1\left(x\cdot\lfloor x\rfloor\right) dx + \displaystyle\int\limits_1^2\left(x\cdot\lfloor x\rfloor\right) dx  $\\

$ \displaystyle\int\limits_{-1}^0\left(x\cdot\lfloor x\rfloor\right) dx  = -\displaystyle\int\limits_{-1}^0xdx =-\left[\dfrac{x^2}{2}\right]_{-1}^0 = \dfrac{1}{2} $\\

$ \displaystyle\int\limits_0^1\left(x\cdot\lfloor x\rfloor\right) dx = 0 \cdot \displaystyle\int\limits_0^1xdx = 0 $\\

$ \displaystyle\int\limits_1^2\left(x\cdot\lfloor x\rfloor\right) dx = \displaystyle\int\limits_1^2xdx = \left[\dfrac{x^2}{2}\right]_1^2 = \dfrac{3}{2} $\\

$ \Rightarrow \displaystyle\int\limits_{-1}^2\left(x\cdot\lfloor x\rfloor\right) dx = \dfrac{1}{2} + 0 + \dfrac{3}{2} = 2 $\\

\begin{thebibliography}{999}
%\bibitem{a} Referentie
\end{thebibliography}
\end{document} 
